\section{\textbf{\texttt{sudo()}} Superuser Access}

\begin{frame}[fragile]{What Is \texttt{sudo()}?}
	\begin{itemize}
		\item \textbf{\texttt{sudo()}} returns the same records in a superuser environment.
		\item Superuser bypasses access rights and record rules.
\begin{block}{Method Syntax}
\begin{lstlisting}
records.sudo()
records.sudo(flag=True)
records.sudo(False)  # leave sudo mode
\end{lstlisting}
\end{block}
	\end{itemize}
\end{frame}

\begin{frame}[fragile]{When to Use \texttt{sudo()}}
	\begin{itemize}
		\item System operations that normal users must trigger but cannot fully access.
		\item Creating technical records (followers, logs, queues) in controlled code.
		\item Cron / server actions that need broader access.
	\end{itemize}
\begin{lstlisting}
# Controlled technical write
self.sudo().write({'system_note': 'Processed by cron'})
\end{lstlisting}
\end{frame}

\begin{frame}[fragile]{When NOT to Use \texttt{sudo()}}
	\begin{itemize}
		\item Do not use \texttt{sudo()} just to ``make the error go away''.
		\item Do not expose sudo results directly to untrusted UI input.
		\item Prefer correct access rights and record rules when possible.
	\end{itemize}
\begin{lstlisting}
# Dangerous pattern
self.env['res.partner'].sudo().search([]).unlink()
\end{lstlisting}
	\begin{itemize}
		\item Ask: does this really need superuser power?
	\end{itemize}
\end{frame}

\begin{frame}[fragile]{\texttt{sudo()} vs \texttt{with\_user()}}
	\begin{itemize}
		\item \texttt{sudo()}: bypass security (superuser).
		\item \texttt{with\_user(user)}: act as that user, still under that user's rights.
	\end{itemize}
\begin{lstlisting}
# Bypass rights
self.sudo().write({'state': 'done'})

# Act as teacher user, respecting rights
self.with_user(teacher).write({'state': 'done'})
\end{lstlisting}
\end{frame}
